% T31 - Télécommunications et RF - créée le 2026-11-10 (audit de couverture du référentiel).
% Thème du référentiel quasi absent des sujets 2013-2026 : présence = MENTIONS
% (fibre optique : E1 2014, E2 2018, E2 appl. 2024 et 2025 ; codage NRZ : E1 2026).
% Une seule question réelle exploitable : E1 2014 Q38 (corrigé officiel E1_Corrige_2014).
% QT2 et QT3 sont des exercices INÉDITS, marqués comme tels.
% Bloc 4 : aucun commentaire de jury sur ces notions dans les rapports 2018-2026
% (recherche fibre/antenne/ligne/modulation/dBm le 2026-11-10) : pas d'encadré de citation.
\expandafter\ifx\csname CAPETINCLUDEMODE\endcsname\relax
\documentclass[11pt,a4paper]{article}
% preamble.tex — packages et commandes communes a toutes les fiches CAPET SII IE
\input{glyphtounicode}
\pdfgentounicode=1
\usepackage{cmap}
\usepackage[T1]{fontenc}
\usepackage[utf8]{inputenc}
\usepackage[french]{babel}
% babel[french] active a \begin{document} des raccourcis de ponctuation (:;!?)
% qui cassent la syntaxe to[...] de circuitikz (erreur "+ or - expected") ;
% on les desactive, ce qui ne change que l'espacement fin avant ces signes.
\AtBeginDocument{\shorthandoff{:;!?}}
\usepackage[a4paper, top=1.6cm, bottom=1.6cm, left=1.4cm, right=1.4cm]{geometry}
\setlength{\headheight}{14.5pt}
\usepackage{amsmath, amssymb, amsfonts}
\usepackage{siunitx}
\sisetup{locale=FR, output-decimal-marker={,}, inter-unit-product=\ensuremath{{}\cdot{}},
         group-minimum-digits=4}
% Unites non fournies par siunitx, declarees ici une fois pour toutes (fiches + exercices).
% Les gardes \ifdefined evitent tout conflit avec une redeclaration locale restee dans un
% fichier : ne pas retirer les gardes, ne pas redeclarer ces unites dans les fichiers.
\DeclareSIUnit{\jour}{j}
\ifdefined\var\else\DeclareSIUnit{\var}{var}\fi   % puissance reactive
\ifdefined\vad\else\DeclareSIUnit{\vad}{VAD}\fi   % puissance deformante
\ifdefined\tr\else\DeclareSIUnit{\tr}{tr}\fi      % tour (tr/min)
\usepackage{array, booktabs, multirow, tabularx}
\usepackage{xcolor}
\usepackage[most, breakable]{tcolorbox}
\usepackage{multicol}
\usepackage{enumitem}
\usepackage{mathtools}
\usepackage{lmodern}
% Refonte visuelle 2026 (direction "moderne coloré") : Lato pour le texte,
% maths sans-serif assorties via newtxsf (charge apres amsmath/amssymb).
\usepackage[default]{lato}
\usepackage{newtxsf}
\usepackage{microtype}
\usepackage{fancyhdr}
\usepackage{tikz}
\usetikzlibrary{positioning,arrows.meta}
\usepackage{circuitikz}
\usepackage{lastpage}

% ── Couleurs par domaine (palette modernisée 2026) ────────────────────
\definecolor{colElec}{RGB}{29, 78, 216}        % Électrotechnique — bleu vif
\definecolor{colElecLight}{RGB}{219, 234, 254} % fond bleu clair
\definecolor{colAuto}{RGB}{21, 128, 61}        % Automatique — vert
\definecolor{colMeca}{RGB}{194, 65, 12}        % Mécanique — orange
\definecolor{colTherm}{RGB}{185, 28, 28}       % Thermique — rouge
\definecolor{colQT}{RGB}{232, 244, 233}        % (hérité, plus utilisé par questiontype)
\definecolor{colJury}{RGB}{255, 248, 225}      % fond jury (crème)
\definecolor{colJuryFrame}{RGB}{202, 138, 4}   % bordure/bandeau jury (ambre)
\definecolor{colGrayLight}{RGB}{245, 245, 247} % fond gris neutre
\definecolor{colSeparator}{RGB}{180, 180, 180} % gris séparateur
\definecolor{colExo}{RGB}{8, 145, 178}          % Banque d'exercices — cyan (distinct des fiches)
% ── Encre et gris de la refonte ───────────────────────────────────────
\definecolor{colInk}{RGB}{23, 32, 42}          % texte courant (quasi-noir)
\definecolor{colSlate}{RGB}{83, 98, 118}       % gris ardoise (méta, en-têtes)
\definecolor{colRule}{RGB}{203, 213, 225}      % filets clairs
\definecolor{colJuryBg}{RGB}{255, 248, 225}    % fond crème jury
\definecolor{colPresence}{RGB}{83, 98, 118}    % bloc 5 (présence dans les sujets)
\AtBeginDocument{\color{colInk}}

% ── Style encadré exercice + corrigé (banque d'exercices) ─────────────
\tcbset{
  exoBox/.style={
    breakable,
    enhanced,
    colback=colExo!8!white,
    colframe=colExo,
    colbacktitle=colExo,
    coltitle=white,
    fonttitle=\bfseries\small,
    left=4pt, right=4pt, top=3pt, bottom=3pt,
    before skip=5pt, after skip=5pt,
  }
}

% ── Bascule enonce / corrige (banque d'exercices) ─────────────────────
% \CAPETCORRIGE est defini en ligne de commande (pdflatex "\def\CAPETCORRIGE{1}\input{...}")
% pour compiler la variante corriges d'un meme fichier source ; par defaut (non defini),
% la compilation produit la variante enonces. \CAPETSHOW{enonce}{corrige} choisit la branche.
%
% TROIS MODES (refonte 2026-09-20, manuel de revision) :
%   - aucun drapeau   -> imprime l'enonce            (variante enonces, INCHANGE)
%   - \CAPETCORRIGE   -> imprime le corrige          (variante corriges, INCHANGE)
%   - \CAPETMANUEL    -> imprime l'enonce ET met le corrige de cote pour l'annexe
%
% Le mode manuel stocke chaque corrige dans une macro (\csgdef), sans ecriture
% fichier ni re-tokenisation : le contenu (circuitikz, accents, %) traverse intact.
% L'annexe les rejoue ensuite dans l'ordre via \CAPETANNEXE.
\usepackage{etoolbox}

\newcounter{CAPETexo}          % numero continu d'exercice dans le manuel
\setcounter{CAPETexo}{0}
\newcounter{CAPETannexe}       % compteur de relecture en annexe

% Titre courant de la notion en cours (pose par \CAPETNOTION), sert d'intertitre annexe
\newcommand{\CAPETnotioncode}{}
\newcommand{\CAPETnotiontitre}{}

% \CAPETNOTION{code}{titre} : ouvre une notion dans le manuel (memorise le libelle,
% et marque le rang du premier exercice de la notion pour l'intertitre d'annexe).
\newcommand{\CAPETNOTION}[2]{%
  \renewcommand{\CAPETnotioncode}{#1}%
  \renewcommand{\CAPETnotiontitre}{#2}%
  \csgdef{CAPETnotionstart\the\numexpr\value{CAPETexo}+1\relax}{#1 -- #2}%
}

% ── \degree robuste ───────────────────────────────────────────────────
% Certains fichiers d'exercices utilisent \degree. Ni le preambule ni siunitx
% ne le definissent, et \providecommand local ne suffit pas : sous hyperref,
% \degree existe deja comme commande factice (\ERROR), donc \providecommand
% s'abstient et l'appel echoue a la compilation du manuel (erreur invisible
% en compilation isolee). On le (re)definit donc inconditionnellement ici,
% apres le chargement des packages, pour tout le corpus.
\AtBeginDocument{\providecommand{\degree}{}\renewcommand{\degree}{\ensuremath{^\circ}}}

% Libelle decoratif : les fichiers d'exercices utilisent aussi \CAPETSHOW dans
% leur \fancyhead et leur bloc de titre, pour afficher "(enonces)" ou "(corriges)".
% Ces appels-la ne sont PAS des exercices : en mode manuel ils ne doivent ni
% consommer un numero, ni alimenter l'annexe. On les reconnait a leur premier
% argument, strictement egal au libelle ci-dessous.
\newcommand{\CAPETLIBELLEENONCES}{(\'enonc\'es)}

\ifdefined\CAPETMANUEL
  % --- Mode manuel : enonce en place, corrige differe vers l'annexe ---
  \newcommand{\CAPETSHOW}[2]{%
    \def\CAPETargun{#1}%
    \ifx\CAPETargun\CAPETLIBELLEENONCES
      #1%            appel decoratif : rendu tel quel, hors numerotation
    \else
      \CAPETSHOWEXO{#1}{#2}%
    \fi
  }
  \newcommand{\CAPETSHOWEXO}[2]{%
    \stepcounter{CAPETexo}%
    \csgdef{CAPETcorrige\theCAPETexo}{#2}%
    \edef\CAPETtmpnum{\theCAPETexo}%
    #1%
    \par\vspace{2pt}%
    {\footnotesize\itshape\color{colExo}%
     \hyperref[capet:cor:\CAPETtmpnum]{Corrige de l'exercice \CAPETtmpnum\ $\rightarrow$ p.~\pageref{capet:cor:\CAPETtmpnum}}%
     \label{capet:eno:\CAPETtmpnum}}%
  }
  % \CAPETANNEXE : rejoue tous les corriges captures, dans l'ordre, avec intertitres
  \newcommand{\CAPETANNEXE}{%
    \setcounter{CAPETannexe}{0}%
    \loop\ifnum\value{CAPETannexe}<\value{CAPETexo}%
      \stepcounter{CAPETannexe}%
      \edef\CAPETcurnum{\theCAPETannexe}%
      \ifcsname CAPETnotionstart\CAPETcurnum\endcsname
        \par\medskip
        {\normalsize\bfseries\color{colExo}\csuse{CAPETnotionstart\CAPETcurnum}}%
        \par\nobreak\smallskip
      \fi
      \begin{tcolorbox}[exoBox, breakable,
        title={Corrige de l'exercice \CAPETcurnum\
               \hfill\normalfont\footnotesize
               \hyperref[capet:eno:\CAPETcurnum]{enonce p.~\pageref{capet:eno:\CAPETcurnum}}}]%
      \label{capet:cor:\CAPETcurnum}%
      \csuse{CAPETcorrige\CAPETcurnum}%
      \end{tcolorbox}%
    \repeat
  }
\else
  % --- Modes historiques, comportement strictement inchange ---
  \newcommand{\CAPETSHOW}[2]{\ifdefined\CAPETCORRIGE #2\else #1\fi}
  \newcommand{\CAPETANNEXE}{}
\fi

% ── Difficulté 1 à 3 (étoiles pleines, amssymb) ───────────────────────
\newcommand{\difficulte}[1]{%
  \ifcase#1\relax
  \or $\bigstar$\;niveau 1/3%
  \or $\bigstar\bigstar$\;niveau 2/3%
  \or $\bigstar\bigstar\bigstar$\;niveau 3/3%
  \fi
}

% ── Style encadré extrait de document technique (banque d'exercices) ──
\tcbset{
  docTech/.style={
    breakable,
    enhanced,
    colback=colGrayLight,
    colframe=colSeparator,
    fonttitle=\bfseries\small\itshape,
    coltitle=black,
    colbacktitle=colGrayLight,
    left=4pt, right=4pt, top=3pt, bottom=3pt,
    before skip=4pt, after skip=4pt,
  }
}

% ── Style tcolorbox de section principale ────────────────────────────
% Refonte 2026 : bandeau titre coloré arrondi + fond très pâle, sans cadre.
\tcbset{
  ficheSection/.style={
    breakable, enhanced jigsaw,
    frame hidden, boxrule=0pt, arc=1.8mm,
    colback=#1!4!white,
    colbacktitle=#1,
    coltitle=white,
    fonttitle=\sffamily\bfseries,
    left=7pt, right=7pt, top=5pt, bottom=5pt,
    toptitle=3pt, bottomtitle=3pt,
    before skip=10pt, after skip=8pt,
  }
}

% ── Style encadré question type ──────────────────────────────────────
% Carte blanche, bandeau sombre, liseré très clair.
\tcbset{
  questiontype/.style={
    breakable, enhanced jigsaw,
    colback=white,
    frame hidden, boxrule=0pt, arc=1.8mm,
    colbacktitle=colElec!78!black,
    coltitle=white,
    fonttitle=\sffamily\bfseries\small,
    left=7pt, right=7pt, top=4pt, bottom=4pt,
    toptitle=2.5pt, bottomtitle=2.5pt,
    before skip=8pt, after skip=8pt,
    borderline={0.6pt}{0pt}{colElec!12!white},
  }
}

% ── Style encadré jury ───────────────────────────────────────────────
% Bandeau ambre, fond crème, sans cadre.
\tcbset{
  jury/.style={
    breakable, enhanced jigsaw,
    colback=colJuryBg,
    frame hidden, boxrule=0pt, arc=1.8mm,
    colbacktitle=colJuryFrame,
    coltitle=white,
    fonttitle=\sffamily\bfseries,
    left=7pt, right=7pt, top=4pt, bottom=4pt,
    toptitle=3pt, bottomtitle=3pt,
    before skip=10pt, after skip=8pt,
  }
}

% ── En-tête / pied de page communs (refonte 2026) ─────────────────────
% #1 code fiche, #2 titre court, #3 couleur domaine
\newcommand{\ficheHeaders}[3]{%
  \fancyhead[L]{\small\sffamily{\bfseries\textcolor{#3}{#1}}\hspace{0.7em}\textcolor{colSlate}{#2}}%
  \fancyhead[R]{\small\sffamily\textcolor{colSlate}{CAPET SII IE\, --- \,Session 2026}}%
  \fancyfoot[C]{\small\sffamily\textcolor{colSlate}{\thepage\,/\,\pageref{LastPage}}}%
  \fancyhead[C]{}\fancyfoot[L]{}\fancyfoot[R]{}%
}

% ── Bloc en-tête de fiche (BLOC 0, refonte 2026) ──────────────────────
% #1 code  #2 titre  #3 couleur domaine  #4 domaine  #5 ligne présence  #6 sous-thèmes
\newcommand{\ficheEntete}[6]{%
  \begin{tcolorbox}[enhanced, breakable,
    frame hidden, boxrule=0pt, arc=2.5mm,
    interior style={left color=#3, right color=#3!75!black},
    left=9pt, right=9pt, top=7pt, bottom=7pt,
    before skip=0pt, after skip=6pt]
    {\color{white}\sffamily\bfseries\LARGE #1\quad #2}\\[4pt]
    {\color{#3!22!white}\sffamily\small \MakeUppercase{#4}\quad\textbar\quad #5}\\[1pt]
    {\color{#3!35!white}\sffamily\footnotesize Sous-thèmes : #6}
  \end{tcolorbox}%
}

% ── Mise en page ─────────────────────────────────────────────────────
\pagestyle{fancy}
\fancyhf{}
\renewcommand{\headrulewidth}{0.5pt}
\renewcommand{\footrulewidth}{0pt}
\AtBeginDocument{%
  \renewcommand{\headrule}{{\color{colRule}\hrule width\headwidth height 0.5pt}}%
}

\setlist[itemize]{leftmargin=1.2em, itemsep=1pt, topsep=2pt, parsep=0pt}
\setlist[enumerate]{leftmargin=1.2em, itemsep=3pt, topsep=2pt, parsep=0pt}

\setlength{\parskip}{2pt}
\setlength{\columnsep}{10pt}

\def\CAPETBEGINFILE{\begin{document}}
\def\CAPETENDFILE{\end{document}}
\else
\def\CAPETBEGINFILE{}
\def\CAPETENDFILE{}
\fi

\ficheHeaders{T31}{Télécoms et RF}{colElec}

\CAPETBEGINFILE

% -----------------------------------------------------------------------
%  BLOC EN-TETE
% -----------------------------------------------------------------------

\ficheEntete{T31}{Télécommunications et RF - Transmission, lignes, antennes}{colElec}%
  {électronique - Télécommunications et radiofréquences}%
  {\textbf{Présence :} \textbf{4/13} sujets (2013-2026) - thème \emph{rare}, au programme}%
  {bande de base NRZ/Manchester - modulations ASK/FSK/PSK/QAM - multiplexage - PLL - lignes et adaptation - fibre optique - antennes et bilan de liaison - exposition aux champs - instrumentation RF}

% -----------------------------------------------------------------------
%  1. RAPPELS DE COURS
% -----------------------------------------------------------------------

\begin{tcolorbox}[
  ficheSection=colElec,
  title={1. Rappels de cours},
]

\textbf{1.1 Bande de base et modulations numériques}

\smallskip
En \textbf{bande de base}, les bits sont émis directement sous forme de niveaux. Un \textbf{symbole} de durée $T_s$ porte $\log_2 M$ bits ($M$ états) : la rapidité de modulation $R = 1/T_s$ s'exprime en \textbf{bauds}, le débit $D = R\log_2 M$ en bit/s (voir T25). Sur un canal idéal de bande $B$, on ne peut pas dépasser $R_{max} = 2B$ (critère de Nyquist).

\begin{center}
\renewcommand{\arraystretch}{1.35}
\begin{tabular}{@{}lll@{}}
\toprule
\textbf{Code} & \textbf{Principe} & \textbf{Fondamental maximal} \\
\midrule
NRZ & niveau haut = 1, niveau bas = 0 & $D/2$ (suite 0101) ; pas d'horloge transmise \\
Manchester & transition au milieu de chaque bit & $D$ ; horloge récupérable, bande doublée \\
\bottomrule
\end{tabular}
\end{center}

Pour transmettre sur un canal passe-bande, on module une porteuse :
\begin{center}
\renewcommand{\arraystretch}{1.35}
\begin{tabular}{@{}llll@{}}
\toprule
\textbf{Modulation} & \textbf{Grandeur modulée} & \textbf{États} & \textbf{Usage} \\
\midrule
ASK / OOK & amplitude & 2 & télécommandes, RFID \\
FSK & fréquence & 2 ou 4 & LoRa en mode FSK, modems (T25) \\
BPSK / QPSK & phase & 2 / 4 & satellites, Wi-Fi (bas débit) \\
QAM-16, 64\ldots & amplitude et phase & 16, 64\ldots & Wi-Fi, 4G/5G, TNT \\
\bottomrule
\end{tabular}
\end{center}

\begin{center}
\begin{tikzpicture}[scale=1.0, font=\small]
  \draw[->, thick] (-1.8,0) -- (1.9,0) node[right] {$I$};
  \draw[->, thick] (0,-1.8) -- (0,1.9) node[above] {$Q$};
  \draw[colElec!50, dashed] (0,0) circle (1.2);
  \foreach \x/\y/\b in {0.85/0.85/00, -0.85/0.85/01, -0.85/-0.85/11, 0.85/-0.85/10} {
    \fill[colElec] (\x,\y) circle (2.2pt);
    \node[colElec!70!black] at (1.45*\x,1.45*\y) {\footnotesize\b};
  }
  \node[align=left, font=\footnotesize] at (3.6,0.9) {QPSK : 4 phases\\$\log_2 4 = 2$ bits/symbole};
  \node[align=left, font=\footnotesize] at (3.6,-0.6) {codage de Gray :\\voisins = 1 bit d'écart};
\end{tikzpicture}
\end{center}

\medskip\hrule\vspace{4pt}

\textbf{1.2 Multiplexage et synthèse de fréquence}

\smallskip
\textbf{Multiplexage fréquentiel (FDM)} : $N$ voies de bande $B$ séparées par des bandes de garde $B_g$ occupent $B_{tot} = N B + (N-1) B_g$. Le multiplexage temporel (TDM) partage au contraire le temps (trames, intervalles).

\textbf{Boucle à verrouillage de phase (PLL)} : un comparateur de phase asservit un oscillateur commandé en tension (VCO) pour que $f_{out}/N = f_{ref}/R$. On obtient $f_{out} = \frac{N}{R}f_{ref}$, réglable par pas de $f_{ref}/R$, avec la stabilité du \textbf{quartz} de référence (dérive en ppm : $\Delta f = f_0\,\delta\cdot 10^{-6}$).

\medskip\hrule\vspace{4pt}

\textbf{1.3 Propagation guidée : lignes et fibres optiques}

\smallskip
Une liaison est une \textbf{ligne} dès que sa longueur dépasse environ $\lambda/10$ : il faut alors tenir compte de la propagation. Une onde s'y propage à $v = c/\sqrt{\varepsilon_r}$ (environ $2\cdot 10^8$\,m/s dans un coaxial) et la ligne présente une impédance caractéristique $Z_c = \sqrt{L'/C'}$ (50\,$\Omega$ en RF, 75\,$\Omega$ en vidéo).

\begin{center}
\renewcommand{\arraystretch}{1.4}
\begin{tabular}{@{}lll@{}}
\toprule
\textbf{Grandeur} & \textbf{Expression} & \textbf{Adaptation parfaite} \\
\midrule
Coefficient de réflexion & $\Gamma = \dfrac{Z_L - Z_c}{Z_L + Z_c}$ & $\Gamma = 0$ \\
Rapport d'ondes stationnaires & $\mathrm{ROS} = \dfrac{1+|\Gamma|}{1-|\Gamma|}$ & $\mathrm{ROS} = 1$ \\
Affaiblissement de réflexion & $\mathrm{RL} = -20\log|\Gamma|$ & $\mathrm{RL} \to \infty$ \\
Puissance réfléchie & $P_r = |\Gamma|^2 P_i$ & $P_r = 0$ \\
\bottomrule
\end{tabular}
\end{center}

\textbf{Fibre optique} : la lumière est guidée par réflexion totale dans le c\oe ur d'indice $n_1 > n_2$ (gaine). Ouverture numérique $\mathrm{ON} = \sin\theta_{max} = \sqrt{n_1^2 - n_2^2}$ ; vitesse $v = c/n_1$ (environ $2\cdot10^8$\,m/s pour $n_1 \approx 1{,}5$). Fibre \textbf{multimode} : courtes distances, dispersion modale ; \textbf{monomode} : longues distances, hauts débits. Atténuation typique : 0{,}2 à 3\,dB/km. Insensible aux perturbations électromagnétiques et isolante.

\medskip\hrule\vspace{4pt}

\textbf{1.4 Antennes, bilan de liaison, exposition, instrumentation RF}

\smallskip
Longueur d'onde $\lambda = c/f$. Dipôle demi-onde : $L \approx \lambda/2$, gain $2{,}15$\,dBi, impédance $\approx 73\,\Omega$ ; monopole quart d'onde sur plan de masse : $L \approx \lambda/4$. Une antenne se choisit par sa bande, son gain, son diagramme de rayonnement, sa polarisation et son impédance (adaptée à la ligne, ROS faible).

\textbf{Bilan de liaison} en champ lointain (formule de Friis) : $P_r = P_e G_e G_r\left(\frac{\lambda}{4\pi d}\right)^2$, soit en décibels $P_r = P_e + G_e + G_r - A_{EL}$ avec $A_{EL} = 20\log\frac{4\pi d}{\lambda}$. La liaison fonctionne si $P_r$ dépasse la sensibilité du récepteur ; l'écart est la \textbf{marge}.

\textbf{Exposition du public} : densité de puissance $S = \frac{P_e G_e}{4\pi d^2}$, champ $E = \frac{\sqrt{30 P_e G_e}}{d}$, avec $S = E^2/377$. Niveaux de référence (recommandation 1999/519/CE, décret 2002-775) : \SI{28}{V/m} de 10 à 400\,MHz, $1{,}375\sqrt{f_{MHz}}$\,V/m de 400 à 2\,000\,MHz (\SI{41}{V/m} à 900\,MHz), \SI{61}{V/m} au-delà de 2\,GHz. En basse fréquence (50\,Hz), on vérifie le champ magnétique (\SI{100}{\micro\tesla}) et le champ électrique (\SI{5}{kV/m}).

\textbf{Instrumentation RF} : l'\textbf{analyseur de spectre} affiche la puissance (dBm) en fonction de la fréquence (réglages : bande, RBW, niveau de référence) ; l'\textbf{analyseur de réseau vectoriel} mesure les paramètres S ($S_{11} = \Gamma$ en réflexion, $S_{21}$ en transmission) ; l'\textbf{analyseur d'impédance} mesure $\underline{Z}(f)$ d'un composant ou d'une antenne.

\end{tcolorbox}

% -----------------------------------------------------------------------
%  2. FORMULAIRE
% -----------------------------------------------------------------------

\begin{tcolorbox}[
  ficheSection=colElec,
  title={2. Formulaire synthétique},
]
\small
\begin{multicols}{2}
\renewcommand{\arraystretch}{1.45}
\begin{tabular}{@{}p{4.3cm}ll@{}}
\toprule
\textbf{Formule} & \textbf{Usage} & \textbf{U.} \\
\midrule
\multicolumn{3}{@{}l}{\footnotesize\textit{Transmission}} \\
$R = 1/T_s$ & Rapidité & Bd \\
$R_{max} = 2B$ & Nyquist & Bd \\
$f_{NRZ} = D/2$ ; $f_{Manch} = D$ & Fondamental & Hz \\
$B_{tot} = NB + (N-1)B_g$ & FDM & Hz \\
$f_{out} = \frac{N}{R} f_{ref}$ & PLL & Hz \\
$P_{dBm} = 10\log(P/1\,\mathrm{mW})$ & Puissance & dBm \\
\midrule
\multicolumn{3}{@{}l}{\footnotesize\textit{Lignes et fibres}} \\
$v = c/\sqrt{\varepsilon_r}$ ; $v = c/n_1$ & Propagation & m/s \\
$Z_c = \sqrt{L'/C'}$ & Ligne & $\Omega$ \\
$\Gamma = (Z_L - Z_c)/(Z_L + Z_c)$ & Réflexion & - \\
$\mathrm{ROS} = (1+|\Gamma|)/(1-|\Gamma|)$ & Désadaptation & - \\
$\mathrm{ON} = \sqrt{n_1^2 - n_2^2}$ & Fibre & - \\
\bottomrule
\end{tabular}

\columnbreak

\renewcommand{\arraystretch}{1.45}
\begin{tabular}{@{}p{4.3cm}ll@{}}
\toprule
\textbf{Formule} & \textbf{Usage} & \textbf{U.} \\
\midrule
\multicolumn{3}{@{}l}{\footnotesize\textit{Antennes et liaison}} \\
$\lambda = c/f$ & Longueur d'onde & m \\
$P_r = P_e G_e G_r (\lambda/4\pi d)^2$ & Friis & W \\
$A_{EL} = 20\log(4\pi d/\lambda)$ & Espace libre & dB \\
$P_r = P_e + G_e + G_r - A_{EL}$ & Bilan & dBm \\
\midrule
\multicolumn{3}{@{}l}{\footnotesize\textit{Exposition}} \\
$S = P_e G_e/(4\pi d^2)$ & Densité & W/m$^2$ \\
$E = \sqrt{30 P_e G_e}/d$ & Champ lointain & V/m \\
$S = E^2/377$ & Onde plane & W/m$^2$ \\
$E_{lim} = 1{,}375\sqrt{f_{MHz}}$ & 0{,}4--2\,GHz & V/m \\
\bottomrule
\end{tabular}
\end{multicols}
\normalsize
\end{tcolorbox}

% -----------------------------------------------------------------------
%  3. QUESTIONS TYPES
% -----------------------------------------------------------------------

\begin{tcolorbox}[
  ficheSection=colElec,
  title={3. Questions types - une issue d'un sujet réel, deux inédites},
]

\begin{tcolorbox}[questiontype,
  title={QT1 - Fibre optique ou paire torsadée \hfill\normalfont\footnotesize E1 commun 2014 (Q38)}
]
\textbf{Énoncé.}
Un paquet de 256 octets parcourt une liaison de \SI{550}{m}. Avec la fibre optique, le débit vaut \SI{1}{Gbit/s} ; avec la paire torsadée, \SI{2}{Mbit/s}. Le corrigé officiel retient des vitesses de propagation de $3\cdot10^8$\,m/s (fibre) et $2{,}26\cdot10^8$\,m/s (cuivre). Déterminer la durée nécessaire pour recevoir le dernier octet dans chaque cas et comparer.

\tcbline

\textbf{Corrigé.} \textit{(d'après le corrigé officiel E1 2014)}

\textbf{Étape 1 - Nombre de bits.} $N = 256 \times 8 = 2\,048$ bits.

\textbf{Étape 2 - Fibre.} Durée d'émission plus durée de propagation :
\[
  t = \frac{N}{D} + \frac{d}{v} = \frac{2\,048}{10^9} + \frac{550}{3\cdot10^8}
    = 2{,}05 + 1{,}83 = \mathbf{3{,}88\ \mu s}
\]

\textbf{Étape 3 - Paire torsadée.}
\[
  t = \frac{2\,048}{2\cdot10^6} + \frac{550}{2{,}26\cdot10^8}
    = 1\,024 + 2{,}43\ \mu\mathrm{s} = \mathbf{1{,}026\ ms}
\]

\textbf{Point de vigilance.} Dans une fibre, la lumière va à $c/n_1 \approx 2\cdot10^8$\,m/s et non à $c$ : la propagation dure alors \SI{2,75}{\micro\second} et le total \SI{4,80}{\micro\second}. La conclusion est inchangée.

\textbf{Conclusion.} La fibre délivre le paquet environ 260 fois plus vite. C'est le débit, donc la durée d'émission, qui fait la différence : la propagation ne pèse que quelques microsecondes dans les deux cas. La fibre est retenue pour respecter le temps de cycle du réseau.
\end{tcolorbox}

\vspace{4pt}

\begin{tcolorbox}[questiontype,
  title={QT2 - Bilan de liaison radio à 868 MHz \hfill\normalfont\footnotesize exercice inédit}
]
\textbf{Énoncé.}
Un capteur LoRa émet $P_e = 14$\,dBm à $f = 868$\,MHz vers une passerelle située à $d = 5$\,km. Les deux antennes sont des dipôles ($G = 2{,}15$\,dBi). La sensibilité du récepteur vaut $-137$\,dBm. (a) Calculer la puissance reçue en espace libre et la marge. (b) Calculer le champ électrique à 1\,m de l'émetteur et le comparer au niveau de référence.

\tcbline

\textbf{Corrigé.}

\textbf{Étape 1 - Longueur d'onde et affaiblissement.}
\[
  \lambda = \frac{3\cdot10^8}{868\cdot10^6} = 0{,}346\ \mathrm{m}
  \qquad
  A_{EL} = 20\log\frac{4\pi\times 5\,000}{0{,}346} = \mathbf{105{,}2\ dB}
\]

\textbf{Étape 2 - Puissance reçue et marge.}
\[
  P_r = 14 + 2{,}15 + 2{,}15 - 105{,}2 = \mathbf{-86{,}9\ dBm}
  \qquad
  \text{marge} = -86{,}9 - (-137) = \mathbf{50\ dB}
\]

\textbf{Étape 3 - Champ à 1 m.} $P_e = 10^{1{,}4}$\,mW $= 25$\,mW et $G = 10^{0{,}215} = 1{,}64$ :
\[
  E = \frac{\sqrt{30 \times 0{,}025 \times 1{,}64}}{1} = \mathbf{1{,}1\ V/m}
  \qquad
  E_{lim} = 1{,}375\sqrt{868} = 40{,}5\ \mathrm{V/m}
\]

\textbf{Conclusion.} En espace libre, la marge de 50\,dB est très confortable ; en milieu urbain, les obstacles en consomment une grande partie, ce qui justifie les facteurs d'étalement élevés de LoRa. Le champ à 1\,m est environ 37 fois inférieur au niveau de référence : l'exposition est conforme.
\end{tcolorbox}

\vspace{4pt}

\begin{tcolorbox}[questiontype,
  title={QT3 - Adaptation d'une antenne sur un coaxial 50 $\Omega$ \hfill\normalfont\footnotesize exercice inédit}
]
\textbf{Énoncé.}
Un câble coaxial ($Z_c = 50\,\Omega$, $\varepsilon_r = 2{,}25$) de \SI{10}{m} alimente à \SI{100}{MHz} une antenne d'impédance $Z_L = 75\,\Omega$. L'émetteur fournit \SI{10}{W}. Faut-il traiter le câble comme une ligne ? Calculer $\Gamma$, le ROS, l'affaiblissement de réflexion et la puissance réfléchie.

\tcbline

\textbf{Corrigé.}

\textbf{Étape 1 - Ligne ou simple liaison ?}
\[
  v = \frac{3\cdot10^8}{\sqrt{2{,}25}} = 2\cdot10^8\ \mathrm{m/s}
  \qquad
  \lambda = \frac{v}{f} = 2\ \mathrm{m}
  \qquad
  \frac{\ell}{\lambda} = 5 \gg \frac{1}{10}
\]
Le câble est une ligne : l'adaptation compte.

\textbf{Étape 2 - Réflexion.}
\[
  \Gamma = \frac{75 - 50}{75 + 50} = \mathbf{0{,}2}
  \qquad
  \mathrm{ROS} = \frac{1{,}2}{0{,}8} = \mathbf{1{,}5}
  \qquad
  \mathrm{RL} = -20\log 0{,}2 = \mathbf{14\ dB}
\]

\textbf{Étape 3 - Puissance réfléchie.} $P_r = |\Gamma|^2 P_i = 0{,}04 \times 10 = \mathbf{0{,}4\ W}$.

\textbf{Conclusion.} 4\,\% de la puissance revient vers l'émetteur ; un ROS de 1{,}5 reste dans la limite usuelle (ROS $\le 2$). Pour l'annuler, on insère un adaptateur d'impédance (transformateur quart d'onde, réseau LC).
\end{tcolorbox}

\end{tcolorbox}

% -----------------------------------------------------------------------
%  4. ATTENTES DU JURY
% -----------------------------------------------------------------------

\begin{tcolorbox}[
  jury,
  title={4. Points de vigilance - erreurs classiques},
]

\begin{enumerate}[label=\textbf{\arabic*.}, leftmargin=*, itemsep=5pt]
  \item \textbf{Confondre dB et dBm.} Le dBm est une puissance absolue (référence 1\,mW), le dB un rapport. On ajoute des dB à des dBm ; additionner deux dBm n'a pas de sens.
  \item \textbf{Confondre bauds et bit/s.} Les deux ne coïncident que pour $M = 2$ états ; en QPSK, 1\,000 bauds font 2\,000 bit/s.
  \item \textbf{Utiliser $c$ dans un câble ou une fibre.} La vitesse vaut $c/\sqrt{\varepsilon_r}$ ou $c/n_1$, soit environ $2\cdot10^8$\,m/s.
  \item \textbf{Oublier la durée d'émission.} La durée d'une transmission additionne $N/D$ (émission des bits) et $d/v$ (propagation) ; c'est souvent la première qui domine.
  \item \textbf{Appliquer Friis en champ proche.} La formule suppose $d \gg \lambda$ et l'espace libre ; en intérieur, la marge doit être nettement plus grande.
  \item \textbf{Comparer un champ à la mauvaise limite.} Le niveau de référence dépend de la fréquence (28, 41, 58 ou 61\,V/m).
  \item \textbf{Négliger l'adaptation d'une ligne longue.} Au-delà de $\lambda/10$, une charge différente de $Z_c$ crée réflexions et ondes stationnaires.
\end{enumerate}

\medskip
{\footnotesize\itshape Aucun rapport de jury 2018--2026 ne commente ces notions : les points ci-dessus sont des erreurs de méthode classiques, non des citations.}

\end{tcolorbox}

% -----------------------------------------------------------------------
%  5. PRÉSENCE
% -----------------------------------------------------------------------

\begin{tcolorbox}[
  ficheSection=colPresence,
  title={Présence dans les sujets analysés (2013-2026)},
]
\begin{center}
\footnotesize\renewcommand{\arraystretch}{1.15}
\begin{tabular}{lccccccccccccc}
\toprule
\textbf{Année} & 13 & 14 & 15 & 16 & 17 & 18 & 19 & 20 & 21 & 22 & 23 & 24 & 25 \\
\midrule
Épreuve commune & - & $\checkmark$ & - & - & - & - & - & - & - & - & - & $\checkmark$ & $\checkmark$ \\
E2 discipl. IE  & - & - & - & - & - & $\checkmark$ & - & - & - & - & - & - & - \\
\bottomrule
\end{tabular}
\end{center}
\smallskip
\footnotesize\textbf{Sous-thèmes rencontrés :} fibre optique contre cuivre (2014, question 38 ; simples mentions en 2018, 2024, 2025) - codage NRZ (E1 2026, question D3). Les autres notions (antennes, lignes, PLL, exposition, instrumentation RF) sont au programme mais n'ont jamais été interrogées de 2013 à 2026.
\end{tcolorbox}

\CAPETENDFILE
